\section{Geometric and lattice comparisons with independent proofs}
\label{sec:sr-geometric}
Recent translations by Baelson, reviewed by lowrank, identify the
principalisation, arithmetic intersection, toroidal-chart and Fock-model
branches of the original development. A statement that merely re-exports
an earlier theorem does not yet explain that theorem's proof.
We record geometric and lattice comparisons that can be proved directly here,
and specify their domains to avoid applying them beyond their hypotheses.

\subsection{An exact comparison of fractional-ideal lattices}
\label{sec:sr-ideal-lattices}
One ingredient of the high-twist class-counting argument is a
determinant-prescribed comparison of integral lattices.
Here it has a direct proof over a Dedekind domain.

\begin{srproposition}\label{prop:sr-ideal-lattices}
Let \(R\) be a Dedekind domain with fraction field \(F\).
Let \(I_1,\ldots,I_r\) and \(J_1,\ldots,J_r\) be nonzero
fractional ideals, and let \(\beta\in F^\times\), with
\[
                  \prod_{j=1}^rJ_j=\beta\prod_{j=1}^rI_j.
\]
If \(r=0\), assume also \(\beta=1\).
There exists \(A\in\mathrm{GL}_r(F)\) with
\[
 \det A=\beta,\qquad
             A\left(\bigoplus_{j=1}^rI_j\right)
                         =\bigoplus_{j=1}^rJ_j.
\]
\end{srproposition}
\begin{proof}
We first construct, for any two nonzero fractional ideals \(I,J\),
a determinant-one matrix sending \(I\oplus J\) onto \(R\oplus IJ\).
Choose \(0\ne a\in I^{-1}\).
Then \(aI\) is a nonzero integral ideal; let \(\mathcal P\) be its
finite set of prime divisors.

Choose \(b\in J^{-1}\) whose image in \(J^{-1}/PJ^{-1}\) is nonzero
for every \(P\in\mathcal P\).
Here is a direct justification of that simultaneous choice.
For any such \(P\), \(J^{-1}\ne PJ^{-1}\), since multiplication
by the invertible ideal \(J\) would otherwise give \(R=P\).
Choose \(b_P\) outside \(PJ^{-1}\).
For distinct maximal ideals the Chinese remainder construction gives
\(e_P\in R\) with
\(e_P\equiv1\pmod P\) and \(e_P\equiv0\pmod Q\) for every other
\(Q\in\mathcal P\).
Then \(b=\sum_P e_Pb_P\) has the required residues.
The Chinese remainder construction used here follows inductively
from \(P+Q=R\) for distinct maximal ideals.
If \(\mathcal P\) is empty, take \(b=0\).

Both \(aI\) and \(bJ\) are integral ideals and
\[
                         aI+bJ=R.
\]
Indeed a maximal ideal containing their sum would contain \(aI\),
so would belong to \(\mathcal P\); but \(bJ\subseteq P\) would
imply \(b\in PJ^{-1}\), contrary to the choice.
Choose \(x\in I,y\in J\) with \(ax+by=1\), and put
\[
       B=\begin{pmatrix}a&b\\-y&x\end{pmatrix},
       \qquad
       B^{-1}=\begin{pmatrix}x&-b\\y&a\end{pmatrix}.
\]
The displayed relation gives \(\det B=1\) and verifies the inverse.
For \(u\in I,v\in J\), the first coordinate \(au+bv\) belongs
to \(R\), and the second coordinate \(-yu+xv\) belongs to \(IJ\).
Conversely, for \(z\in R,w\in IJ\),
\[
                 xz-bw\in I,\qquad yz+aw\in J,
\]
using \(b\in J^{-1},a\in I^{-1}\).
Thus \(B\) induces the claimed exact lattice comparison.

For \(r\geq2\), apply these two-coordinate matrices successively
in coordinates \((1,2),(2,3),\ldots,(r-1,r)\).
All full matrices have determinant one, and their product \(B_I\)
sends
\[
       \bigoplus_j I_j\quad\hbox{onto}\quad
                             R^{r-1}\oplus\prod_j I_j .
\]
Construct \(B_J\) in the same manner.
The diagonal matrix
\(D=\operatorname{diag}(1,\ldots,1,\beta)\) sends the first
normalized lattice onto the second and has determinant \(\beta\).
Therefore \(A=B_J^{-1}DB_I\) has the required determinant and
image. For \(r=1\), take \(A=(\beta)\).
For \(r=0,\beta=1\), take the identity of the zero space,
whose determinant is one.
\end{proof}

In a number field \(F\), write
\[
 I_j=\prod_P P^{a_{j,P}},\qquad
 J_j=\prod_P P^{b_{j,P}},
\]
with integer exponents of finite support.
If
\[
              v_P(\beta)=\sum_j(b_{j,P}-a_{j,P})
                                      \quad\hbox{for every }P ,
\]
unique ideal factorization gives
\(\prod J_j=\beta\prod I_j\).
The proposition consequently produces the exact equivalence
\begin{equation}\label{eq:sr-ideal-valuations}
 \bigl(v_P(x_j)\geq a_{j,P}\ \hbox{for all }j,P\bigr)
 \quad\Longleftrightarrow\quad
 \bigl(v_P((Ax)_j)\geq b_{j,P}\ \hbox{for all }j,P\bigr).
\end{equation}
We use \(v_P(0)=+\infty\).
Membership in a fractional ideal is equivalent to these valuation
inequalities, by unique factorization and localization at the
discrete valuation rings \(R_P\).
This proves, in particular, the determinant and valuation comparison
called \srcid{SStein} in the class-counting chain leading to \(K_2\).

There is also an exact archimedean consequence.
Use the standard Minkowski embedding of \(F^r\) into
\(\mathbb R^{rr_1}\times\mathbb C^{rr_2}\), with fixed Lebesgue
measure, and let \(\operatorname{covol}\) denote lattice covolume.
The two direct sums of fractional ideals are full lattices, and
\begin{equation}\label{eq:sr-ideal-covolume}
 \operatorname{covol}\left(\bigoplus_jJ_j\right)
   =|\operatorname{Norm}_{F/\mathbb Q}\beta|\,
                  \operatorname{covol}\left(\bigoplus_jI_j\right).
\end{equation}
Indeed the real absolute Jacobian of \(A\) is
\[
 \prod_{\sigma\ {\rm real}}|\sigma(\det A)|
 \prod_{\sigma\ {\rm complex}/\mathrm{conj}}|\sigma(\det A)|^2
                 =|\operatorname{Norm}_{F/\mathbb Q}\beta|.
\]
A complex-linear matrix has real determinant equal to the square
of the absolute value of its complex determinant; this follows
by triangularization, or by writing its real and imaginary blocks.
Applying the change-of-variables formula to a lattice fundamental
parallelepiped proves \eqref{eq:sr-ideal-covolume}.
This comparison applies uniformly at every layer and twist.
The later metric estimates also require bounds on operator norms;
the exact determinant comparison supplies the covolume factor.

\subsection{Logarithmic norms after principalisation}
Let \(X\) be a connected smooth projective complex variety of dimension \(d\),
\(L\) a holomorphic line bundle with smooth positive Hermitian norm,
\(s\) a nonzero holomorphic section, and \(\alpha\) a smooth top-degree
real form on \(X\).
Then
\begin{equation}\label{eq:sr-logint}
 \int_X|\log\|s\||\,|\alpha|<\infty.
\end{equation}
In particular this applies when \(\alpha\) is a product of smooth Chern
forms of total degree \(d\).

Here is the proof, with principalisation as the classical resolution
theorem in characteristic zero. Take a proper modification
\(\pi:Y\to X\) with \(Y\) smooth, such that the zero divisor of
\(\pi^*s\) is a simple normal crossings divisor.
On a coordinate polydisk of \(Y\), with a nonvanishing local frame,
\(\pi^*s=u z_1^{a_1}\cdots z_k^{a_k}\), where \(u\) is holomorphic
and nowhere zero and the \(a_i\) are nonnegative integers.
On a slightly smaller closed polydisk, both the norm of the frame
and \(|u|\) are bounded above and below by positive constants.
Consequently
\[
 |\log\|\pi^*s\||\leq C+\sum_{i=1}^ka_i|\log|z_i||.
\]
The smooth form \(\pi^*\alpha\) has bounded coefficients there.
Fubini and polar coordinates give
\[
 \int_{|z|\leq r}|\log|z||\,dA(z)
       =2\pi\int_0^r\rho|\log\rho|\,d\rho<\infty.
\]
Thus the required absolute integral is finite on each such polydisk.
Compactness supplies a finite cover and a partition of unity.
The modification is a biholomorphism away from analytic sets of measure
zero. Change of variables there transfers the finite absolute integral
to \(X\), proving \eqref{eq:sr-logint}.

This proof covers smooth metrics. It does not assert the corresponding
uniform estimates for every singular cusp metric appearing in the
arithmetic construction. Bounds at those cusps must be checked with
their actual metric growth and boundary multiplicities.

\subsection{The shifted positive-definite lattice identity}
Let \(V\) be a \(d\)-dimensional Euclidean space and \(L\subset V\)
a full lattice of covolume \(v\). Define
\(L^*=\{\lambda:\langle\lambda,\ell\rangle\in\mathbb Z
\text{ for all }\ell\in L\}\).
For \(t>0,u\in V\),
\begin{equation}\label{eq:sr-latticetheta}
 \sum_{\ell\in L}e^{-\pi t\|\ell+u\|^2}
 =\frac{t^{-d/2}}v
   \sum_{\lambda\in L^*}
     e^{-\pi\|\lambda\|^2/t}e^{2\pi i\langle\lambda,u\rangle}.
\end{equation}
Both sides converge absolutely, locally uniformly in \(t>0,u\);
all derivatives in \(u\) and any fixed number in \(t\) do as well.

To prove this, choose a basis of \(L\) and write it as \(A\mathbb Z^d\).
The least singular value of \(A\) is positive, so
\(\|A n+u\|\geq c\|n\|-\|u\|\).
A Gaussian in \(\|n\|\) dominates the summands and every fixed derivative
on compact parameter sets; sums of these Gaussians converge by comparison
with the product of \(d\) one-dimensional Gaussian sums.
The periodization \(F(u)=\sum_{\ell\in L}e^{-\pi t\|\ell+u\|^2}\)
is smooth and \(L\)-periodic. Its Fourier coefficient at \(\lambda\in L^*\)
is, by integration over a fundamental domain and then all its translates,
\[
 \frac1v\int_Ve^{-\pi t\|x\|^2}e^{-2\pi i\langle\lambda,x\rangle}\,dx
 =\frac{t^{-d/2}}v e^{-\pi\|\lambda\|^2/t}.
\]
The integral is a product of the one-dimensional Gaussian transforms
proved in Section~\ref{sec:sr-transform}.
The coefficients form an absolutely convergent Fourier series, whose
sum therefore equals \(F\). This proves the identity and the derivative
claims.

The factor \(v\) is fixed by the Euclidean measure used in that integral.
If \(L'\subset L\) has finite index \(h\), then
\(\operatorname{covol}(L')=h\operatorname{covol}(L)\):
a fundamental domain for \(L'\) is the disjoint union, up to its boundary,
of \(h\) translates of a fundamental domain for \(L\).
For an invertible real linear map \(B\),
\(\operatorname{covol}(BL)=|\det B|\operatorname{covol}(L)\),
by the ordinary change-of-variables formula.
Thus no unspecified volume normalization is hidden in
\eqref{eq:sr-latticetheta}.

Positive definiteness is essential to this particular proof.
For an indefinite quadratic form the displayed unsmoothed Gaussian
need not even tend to zero along the lattice.
Nor does \eqref{eq:sr-latticetheta} establish the required matching of
ramified local Weil sections. The relevant Schwartz functions,
regularizations and local measures in the global theta construction
must be supplied separately.
